\documentclass[12pt,oneside]{article}
\usepackage{etoolbox}
\usepackage[T1]{fontenc}
\usepackage[full]{textcomp}
\usepackage{fancyhdr}
\usepackage{longtable}
%\usepackage[dvips]{graphics,color}      % usual driver for dvi or ps viewer
\usepackage[]{graphicx,color}      % usual driver for converting to pdf later
\usepackage[bottom=1in]{geometry}
\usepackage{verbatim}
%\usepackage[cmintegrals,cmbraces]{newtxmath}
\usepackage[intlimits]{amsmath}    % need for subequations
%\usepackage[]{amsmath}    % need for subequations
\usepackage{amsfonts}
\usepackage{amssymb}
%\usepackage{ebgaramond-maths}
\usepackage[sb]{ebgaramond} %semi-bold option
%\input{../ebgaramond-missing.tex}
\usepackage[boxed]{algorithm2e}
\usepackage{tikz}
\usepackage{float}
\usepackage{courier}
\usepackage{listings}
\usepackage{enumitem}
\usepackage{setspace}
\usepackage{caption}
\usepackage{longtable}
\usepackage{cellspace}
%\usepackage{multirow}
%\usepackage{subsection}

%\renewcommand{\rmdefault}{ugm} %old Garamond
\newcommand{\ttt}[1]{\texttt{\footnotesize{#1}}}
\graphicspath{./}
\DeclareGraphicsExtensions{.pdf,.eps}

\newcommand{\beqn}{\begin{equation}}
\newcommand{\eeqn}{\end{equation}}
\newcommand{\beqnn}{\begin{equation*}}
\newcommand{\eeqnn}{\end{equation*}}
\newcommand{\half}{\text{\scriptsize\ensuremath{\frac{1}{2}}}}
\newcommand{\redsize}[1]{\text{\scriptsize\ensuremath{#1}}}
\newcommand{\bsym}[1]{\boldsymbol{#1}}
\newcommand{\tinysize}[1]{\text{\tiny\ensuremath{#1}}}
\newcommand{\eee}{\text{\large{\rm e}}}
\newcommand{\E}{{\rm I\kern-.3em E}}
\newcommand{\pr}{\mathbb{P}}
\newcommand{\var}{{\rm Var}}
\newcommand{\defas}{\overset{*}{=}}
\newcommand{\sL}{{L}}
\newcommand{\edge}[2]{#1\!\!\rightarrow\!\!#2}
\newcommand{\sedge}[2]{#1\rightarrow#2}
\newcommand{\sumno}{\sum\nolimits}
\newcommand{\ns}[1]{\mbox{\normalsize $#1$}}

\let\osetminus\setminus
\renewcommand{\setminus}{\!\osetminus\!}
\let\oin\in

\renewcommand{\in}{\!\oin\!}

\numberwithin{equation}{section}
\renewcommand{\thefootnote}{\fnsymbol{footnote}}
\DeclareMathOperator*{\argmin}{arg\,min}
\DeclareMathOperator*{\argmax}{arg\,max}
\DeclareMathOperator*{\mino}{min}
\DeclareMathOperator*{\maxo}{max}
\DeclareMathOperator*{\erf}{erf}
\DeclareMathOperator*{\arcosh}{arcosh}
\DeclareMathOperator*{\sgn}{sign}
\DeclareMathOperator*{\forallo}{\Large\forall}
\newcommand{\ie}{{\em i.e.}\ }
\newcommand{\eg}{{\em e.g.}\ }
\newtheorem{defn}{Definition}
\newtheorem{lemma}{Lemma}
\newtheorem{cor}{Corollary}
\newtheorem{prop}{Proposition}
\newtheorem{theorem}{Theorem}
\newenvironment{proof}[1][]{\paragraph{Proof #1:}}{\hfill$\square$}
\patchcmd{\thebibliography}{\section*{\refname}}{\section*{\refname}}{}{}
\let\oabsname\abstractname
\renewcommand{\abstractname}{\large{\oabsname}}

\newenvironment{Itemize}{\vspace{-0.5\parskip}\begin{itemize}}{\end{itemize}}
\newenvironment{Enumerate}{\vspace{-0.5\parskip}\begin{enumerate}}{\end{enumerate}}
%\setlist*[itemize]{noitemsep, topsep=0pt}
\setlist*[itemize]{itemsep=5pt, topsep=0pt}
\setlist*[enumerate]{itemsep=5pt, topsep=0pt}

%\definecolor{bggrey}{rgb}{0.95,0.95,0.95}
\definecolor{bggrey}{rgb}{0.93,0.95,0.94}
%\definecolor{bggrey}{rgb}{239,241,242}

\lstdefinestyle{mystyle} {
  framexleftmargin=2pt, %keeps numbers outside of frame
  %framextopmargin=15pt,
  %framexbottommargin=-15pt,
  frame=tb,
  framesep=7pt,
  rulecolor=\color{bggrey},
  backgroundcolor=\color{bggrey},
  basicstyle=\ttfamily\scriptsize,
  linewidth=1.0\textwidth,
  numbers=left,
  showlines=true,
  numberblanklines=false,
  stepnumber=1,
  numbersep=5pt,
  %numberstyle=\ttfamily\scriptsize,
  breaklines=true,
  breakatwhitespace=true,
  tabsize=2,
  keepspaces=true,
  mathescape=true,
  escapeinside={(*@}{@*)}
}

\lstdefinestyle{mystyle1} {
  %backgroundcolor=\color{bggrey},
  %basicstyle=\ttfamily\scriptsize,
  linewidth=0.5\textwidth,
  numbers=none,
  belowskip=-0.8 \baselineskip,
  %breaklines=true,
  %breakatwhitespace=true
}
\lstset{style=mystyle}

% don't need the following. simply use defaults
\setlength{\baselineskip}{16.0pt}    % 16 pt usual spacing between lines
% Default margins are too wide all the way around.  I reset them here
\setlength{\parskip}{2ex plus 2pt}
\setlength{\parindent}{0pt}
\setlength{\oddsidemargin}{.125in}
%\setlength{\oddsidemargin}{0.5cm}
\setlength{\evensidemargin}{0.5cm}
\setlength{\marginparsep}{0.75cm}
\setlength{\marginparwidth}{2.5cm}
\setlength{\marginparpush}{1.0cm}
\setlength{\textwidth}{6.25in}
%\setlength{\textwidth}{150mm}
\setlength{\textheight}{9in}
\setlength{\topmargin}{-.5in}

\renewcommand{\headrulewidth}{0pt}
\renewcommand{\footrulewidth}{0pt}
\fancypagestyle{fpstyle}{ %define a style for fitst page
  \fancyhf{}
  \fancyfoot[L]{{\small\em Working Paper}}
  %\fancyfoot[R]{{\small\em Confidential}}
}
\fancyhf{} %clear all header and footer fields
\fancyfoot[L]{{\small\em Working Paper}}
%\fancyfoot[R]{{\small\em Confidential}}
\fancyfoot[C]{\thepage}
                 
\title{
  \textsc{SABR Simulation for Risk Management}\\%\footnotemark
  \large{\em{(Working Paper)}}
}
\author{Maged Tawfik%\footnotemark %\\
%Citibank
}
\renewcommand{\today}{\vspace{-20pt}\it \small{January 12, 2026}}

\begin{comment}
\end{comment}

\begin{document}
\pagestyle{fancy}

%\leftdisplays
\maketitle
\thispagestyle{fpstyle} % use if do not want page numbers
%\pagestyle{empty} % all pages are not numbered
\begin{center}
\end{center}
\footnotetext[1]{The opinions expressed in the paper are those of the author and do not necessarily reflect the view of Citibank.}
\footnotetext[2]{email: maged.tawfik@citi.com}
%\begin{abstract}
%\input{liqNetflowAbstract_txt.tex}
%\end{abstract}
\section{Introduction}
The SABR (Acronym for Stochastic alpha-beta-rho) model is a stochastic volatility model
first studied by Hagan {\em et al}, 2002\cite{HKLW}.
Due to its ability to replicate certain observed market phenomena, such as
implied volatility smiles, frowns and skews, it has gained wide interest
for the modeling of options.
With its modeling power that can cover underlying assets that display normal, log-normal
or CEV (Constant Elasticity of Variance) behavior,
SABR is currently in extensive use in the area of fixed income derivatives.

Formally, assuming a filtered probability space $(\Omega,\mathcal{F},\mathcal{F}_t,P)$,
SABR is given by the following equations
\begin{subequations}\label{eq:sabr}
    \begin{align}
        dF_t=\,&F_{t}^{\beta} \sigma_t dW_t\quad&\beta\in[0,1],\,F_0>0\label{eq:sabr_a}\\
        d\sigma_t =\,&\nu \sigma_t dZ_t&\nu\ge 0,\,\sigma_0>0\label{eq:sabr_b}\\
        [dZ_t,dW_t] =\,& \rho dt&\rho \in[-1,1]
    \end{align}
\end{subequations}

where $F_t, \sigma_t\in\mathcal{F}_t$, $[.,.]$ is the quadratic variation operator and
$dW_t, dZ_t$ are standard Wiener increments.

Pricing derivatives based on the SABR model is done using either a PDE (Partial Differential Equation)
solver based on the two-dimensional discretization of the Kolmogorov backward equation,
or a Monte Carlo simulation approach.
PDE solvers tend to have much faster convergence.
On the other hand, each run is effective in pricing a single option.
Conversely, a set of Monte Carlo paths, once generated, can be used to price any number of
derivatives as long as they share the same underlying asset and expiry date.
This makes Monte Carlo simulation a more favorable approach in risk management,
where risk measures associated with very large fixed income derivative portfolios need
to be evaluated.

\section{Difficulties in SABR Modeling}
\label{sec:diff}
To properly simulate the SABR model, we need to counter two essential difficulties.
To explain the first difficulty, we use the following discretized version of \eqref{eq:sabr}

\beqn\label{eq:sabr_disc}
\begin{aligned}
    F_t - F_{t-\Delta t} = & \bar{F}_t^\beta \sigma_? W_{\Delta t}\\
    \sigma_t - \sigma_{t-\Delta t} =&\nu \bar{\sigma}_t Z_{\Delta t}\\
    [W_{\Delta t}, Z_{\Delta t}] = \rho \Delta t
\end{aligned}
\eeqn
where $\bar{F}_t$ is an average value of $F_s$ between time $t-\Delta t$ and $t$,
$\bar{\sigma}_t$ is an average value of $\sigma_s$ between time $t-\Delta t$ and $t$,
and $\sigma_?$ is a representative value of $\sigma$ between $t-\Delta t$ and $t$.
The difficulty in choosing an appropriate value $\sigma_?$, when $\nu>0$,
stems from the fact
that the continuous change of the value of $\sigma$ is accompanied by
a continuous change in $F$ dependent on $\sigma$.
Islah, 2009\cite{I}, showed that when $\beta \in \{0,1\}$ or when $\rho=0$,\,
$\sigma_?$ can be exactly represented by the conditional average variance given as
\beqn\label{eq:ith}
\sigma_0\sqrt{\Delta t\;I_{t-\Delta t}^t(\sigma_t)}
\eeqn
where
\beqn
I_{t-\Delta t}^t(\sigma_t):=\frac{1}{\sigma_0^2\, \Delta t}\,
\int_{t-\Delta t}^t \sigma_s^2\, ds \;\bigg|_{\sigma_t}
\eeqn
Choi {\em et al}, 2025\cite{CHK}, give the following equivalent functional
\beqn
I_{t-\Delta t}^t(\hat{Z}_1)\sim \int_0^1 \exp(2\nu\Delta t\, \hat{Z}_s)\,ds\,\big|_{\hat{Z}_1}
\eeqn
where $\hat{Z}_s\sim \mathcal{N}\left(-s\,\nu\sqrt{\Delta t}, \sqrt{s}\right)$.
Although $I_{t-\Delta t}^t(\sigma_t)$ has a closed form PDF, as shown by
Cai {\em et al}, 2017\cite{CSC}, there is no closed form for the CDF and
its numerical integration is very difficult and highly error-prone.
In their paper, Cai {\em et al} suggested the use of Laplace transform inversion
to obtain an accurate and efficient representation of the CDF.
We will be using a modified version of their Laplace inversion algorithm for the
purposes of this paper.

The second difficulty in simulating a SABR process,
is that the proxy for $\sigma_?$ given by \eqref{eq:ith}
is no longer valid.
No other exact proxy is currently available.
As a result, one has to resort to plausible approximations such as those presented in
\cite{CSC},\cite{CHK} and\cite{HKLW}.

\section{Sampling $\sigma_t$}
In the process of simulating the SABR process, one may need to directly sample
the volatility process itself.
As per Equation~\eqref{eq:sabr_b}, $\sigma_t$ follows a log-normal process.
As a result, it can be sampled according to
\beqn\label{eq:sig_sim}
\begin{aligned}
\sigma_t=&\sigma_0\,\exp(\nu\,\sqrt{t}\,\hat{Z}_t)\\
\hat{Z}_t\sim&\mathcal{N}\left(-\half\nu\,\sqrt{t},\,1\right)\;.
\end{aligned}
\eeqn
Thus, $\sigma_t$ is sampled {\em via} a single Gaussian variate.

\section{Sampling the Average Variance}
As explained in Section~\ref{sec:diff}, the average volatility variance $I_0^t(\sigma_t)$
is needed for the simulation of the price process.
As discussed earlier, the sampling of $I_0^t(\sigma_t)$ is fraught with difficulty,
due to the lack of a closed form representation of its CDF.
For the sake of presentational simplicity, we may notationally drop the $\sigma_t$ dependency of
$I_0^h$, in the remainder of the text, without changing the fact.

In\cite{CHK}, Choi {\em et al} proposed an approach for the sampling of $I_0^t$
by matching its CDF to that of shifted log-normal ($\mathcal{SLN}(\mu,\sigma^2,\lambda)$)
process with matching
first four moments.
The $\mathcal{SLN}$ process is defined as
\beqn
\begin{split}
Y\sim&\mathcal{SLN}(\mu,\sigma^2,\lambda) :=\\
Y\sim& \mu\left[(1-\lambda)+
\lambda\exp(\sigma X-\frac{\sigma^2}{2})\right]\quad\quad\mu>0,\sigma>0,0<\lambda\le 1\,, 
\end{split}
\eeqn
where $X$ is a standard Normal variate, where $\mu$ and $\sigma^2$ are the
mean and variance of the LN process and $\lambda$
is the shift parameter.
In their work,
Choi {\em et al} suggested setting $\lambda=\frac{5}{6}$ and sampling $I_0^t$
according to
\beqn\label{eq:choi_i}
\begin{aligned}
I_0^t\sim\mathcal{SLN}\left(\mu,\sigma^2,\frac{5}{6}\right)\\
\sigma = \sqrt{\ln\left(1+\frac{36}{25}v^2\right)}\,,
\end{aligned}
\eeqn
where $\mu$ and $v$ are drived from
\beqn
\begin{aligned}
    \hat{Z}=&\,\frac{\ln(\sigma_t)-\ln(\sigma_0)}{\nu\sqrt{t}}\\
    m_k(\hat{Z}) = &\,\frac{\Phi(\hat{Z}+k\nu\sqrt{t})-\Phi(\hat{Z}-k\nu\sqrt{t})}
    {2k\nu\sqrt{t}\;\phi(\sqrt{\hat{Z}^2+k^2\nu^2t})}\\
    c(\hat{Z})=&\,\half\left(\frac{\sigma_t}{\sigma_0}+\frac{\sigma_0}{\sigma_t}\right)\\
    \mu = &\,\frac{\sigma_t}{\sigma_0}\,m_1(\hat{Z})\\
    v=&\,\left(\frac{\sigma_t}{\sigma_0}\right)
    \frac{\sqrt{m_2(\hat{Z})-c(\hat{Z})\,m_1(\hat{Z}})}{\mu\,\nu\sqrt{t}}
    \,-\mu\,,
\end{aligned}
\eeqn
where $\Phi(.)$ and $\phi(.)$ are standard normal CDF and PDF, respectively.
Significant testing showed that while the approximation approach of \eqref{eq:choi_i}
works reasonably well for $t$ less than three years,
it tends to diverge for larger time steps.

A more robust approach for the sampling of $I_0^t$ is the exact approach of Cai {\em et al}\cite{CSC},
which is based on Laplace transform technique.
In their paper they showed that the Laplace transform $L(\theta)$ of the CDF of $1/I_0^t$ is given by
\beqn\label{eq:laplace}
\begin{aligned}
    L(\theta)=\,&\mathcal{L}(\theta;F_{1/I_0^t}(u))=\int_0^\infty F_{1/I_0^t}(u)\,\exp(-\theta\,u)\,du\\
    L(\theta)=\,&\frac{\sigma_0^2t}{\theta}\,\exp\left\{\frac{[\ln(\sigma_t/\sigma_0)]^2-[\psi(\theta\nu^2/\sigma_0^2)]^2}{2\nu^2\,t}\right\}\\
    \psi(x)=\,&\arcosh\left(\frac{x\sigma_0}{\sigma_t}+\half(\frac{\sigma_t}{\sigma_0}+\frac{\sigma_0}{\sigma_t})\right)\,.
\end{aligned}
\eeqn

Several powerful numerical algorithms have been proposed in the literature for the
inversion of Laplace transforms.
We present here the Euler algorithm proposed by Abate and Whitt (2006)\cite{AW}
for the inversion of \eqref{eq:laplace}, which reads

\beqn\label{eq:euler}
\begin{aligned}
F_{1/I_0^t}(t) =\,&\frac{10^\frac{M}{3}}{t}\,\sum_{k=0}^{2M}\,(-1)^k\eta_k\;\Re\left(L(\frac{\beta_k}{t})\right)\\
\beta_k=\,&\frac{1}{3}\,M \ln(10)\,+\,\pi i k\\
\eta_0=\,&\frac{1}{2},\quad\quad\eta_k=1,\;\; 1\le k\le M\quad\quad\eta_{2M}=&2^{-M}\\
\eta_{2M-k}=\,&\eta_{2M-k+1}\,+\,2^{-M}{M\choose k},\;\; 0<k<M\,,
\end{aligned}
\eeqn
where $i=\sqrt{-1}$, $\Re(.)$ is the real part of a complex number
and $M$ is chosen according to the desired level of accuracy.
The number of correct significant digits is greater than $0.6M$,
for well-behaved transforms.
Notice that the CDF of $I_0^t$ is equal to $1-F_{1/I_0^t}$, since $I_0^t>0$,
and hence, the relationship between the two is a bijection and
\beqnn
\sgn(\frac{d(1/I_0^t)}{dt})=-\sgn(\frac{dI_0^t}{dt})\,.
\eeqnn.

\section{Sampling the CEV Process}\label{sec:cev}
Since the SABR model is a stochastic volatility generalization of the Constant Elasticity of
Variance (CEV) model, sampling CEV variates is usually a needed step in the SABR
simulation.
\begin{defn}[CEV]
A CEV distribution which is denoted by
\beqn
F_t\sim\mathcal{CEV}_\beta(F_0,\sigma^2_0t)\,,
\eeqn
is the distribution of $F_t$, which follows the stochastic process
\beqn
dF_t = \sigma_0 F_t^\beta dW_t\quad \sigma\!\ge\! 0,\, F_0\!\ge\! 0,\, 0\!\le\!\beta\!\le\! 1\,,
\eeqn
where $dW_t$ is a standard Wiener process.
\end{defn}

Islah, 2009\cite{I}, showed that a SABR model with $\rho=0$ is reducible to a CEV model by replacing
$\sigma^2_0t$ by $I_0^t$, and that
the CEV CDF follows the following probability law
\beqn\label{eq:cev}
\begin{aligned}
\pr(F_t=0\big|F_0,\sigma_0,\sigma_t,I_0^t)=&\,1\,-\,Q_{\chi^2}\left(A_0;\frac{1}{1-\beta}\right)\\
\pr(F_t\le u\big|F_0,\sigma_0,\sigma_t,I_0^t)=&\,1\,-\,Q_{\chi\prime^2}\left(A_0;\frac{1}{1-\beta};C_0(u)\right)
\quad u>0\,,
\end{aligned}
\eeqn
where
\beqnn
\begin{aligned}
    A_0=&\,\frac{1}{I_0^t}\left(\frac{F_0^{1-\beta}}{1-\beta}\right)^2\,\text{and}\\
    C_0(u)=&\,\frac{1}{I_0^t}\left(\frac{u^{1-\beta}}{1-\beta}\right)^2\,,
\end{aligned}
\eeqnn
and where $\chi^2(\delta)$ is the chi-squared distribution
with $\delta$ degrees of freedom, and
$\chi\prime^2(\delta,\gamma)$ is the non-central chi-squared distribution
with $\delta$ degrees of freedom and $\gamma$ non-centrality parameter.
The sampling of $F_t$ requires the numerical inversion of Equation
\eqref{eq:cev}, a process fraught with difficulty due
to wide range of the parameters involved which
lead to high degree of numerical instability.

A novel approach proposed by Choi {\em et al},\cite{CHK} based on the
work of Makarov and Glew\cite{MG} and Kang\cite{K}
allows for the direct sampling of a CEV variate through the
drawing of two gamma samples and one Poisson sample.
Using their approach, we define
\beqnn
\begin{aligned}
    z_t:=&\,z(F_t)\\
    z(y):=&\,\frac{u^{2(1-\beta)}}{(1-\beta)^2\sigma_0^2tI_0^t}\,.
\end{aligned}
\eeqnn
A sample of $z_t$ can be drawn according to
\beqn
\begin{aligned}
    Z_t\sim&\, 2\mathcal{G}(N+1)\\
    N\sim&\, \mathcal{POIS}(\frac{z_0}{2} - \,X)\\
    X\sim&\, \mathcal{G}\left(\frac{1}{2(1-\beta)}\right)\\
    F_t=&
    \begin{cases}
        0\quad\quad X\ge 0.5\,z_0\\
        \left(\sigma_0(1-\beta)\sqrt{t\,Z_t\,I_0^t}\right)^\frac{1}{1-\beta}
        \quad \quad X<0.5\,z_0\,,
    \end{cases}
\end{aligned}
\eeqn
where $\mathcal{G}(\alpha)$ is the gamma distribution with shape parameter
$\alpha$ and unit scale parameter, and
$\mathcal{POIS}(\lambda)$ is the Poisson distribution with intensity $\lambda$.

\section{Approximating the SABR Process}
\label{sec:approx}
As already mentioned in Section~\ref{sec:diff}, the SABR process can be
reduced to a CEV process, with $\sigma_0\sqrt{t\,I_0^t}$ playing the
r\^ole of $\sigma_t$, whenever $\rho=0$.
As of the time of this writing, however, there is no known analytic expression
of the general case of SABR dynamics with $\rho\ne 0$.
In this section we present an approximation,
based on an operator splitting technique,
as suggested by Choi {em et al}\cite{CHK}.
According to this approach we rewrite Equation~\eqref{eq:sabr} as
\beqn\label{eq:sabr1}
\frac{dF_t}{F_t^\beta}=\,\sigma_t\,(\rho\,dZ_t \,+\,\sqrt{1-\rho^2}\,dZ_t^\perp)\quad
\text{and}\quad\frac{d\sigma_t}{\sigma_t}=\,\nu\,dZ_t\,,
\eeqn
where $dZ_t$ and $dZ_t^\perp$ are uncorrelated Wiener increments.
Applying the operator splitting technique Equation~\eqref{eq:sabr1} becomes
\begin{subequations}
\begin{align}
\frac{dF_t^{(1)}}{[F_t^{(1)}]^\beta}=&\,\rho\,\sigma_t\,dZ_t\;\quad\text{and}\;\quad
\frac{d\sigma_t}{\sigma_t}=\nu\,dZ_t\label{eq:sabr2a}\\[5pt]
\frac{dF_t^{(2)}}{[F_t^{(2)}]^\beta}=&\,\sqrt{1-\rho^2}\,\sigma_t\,dZ_t^\perp\;\quad\text{and}\quad
\;\frac{d\sigma_t}{\sigma_t}=\nu\,dZ_t\label{eq:sabr2b}\,.
\end{align}
\end{subequations}

The operator splitting technique approximates the solution $F_t$ of \eqref{eq:sabr1},
with initial condition $F_0$,
by $F_t\!\approx\!\! F_t^{(2)}\!\left(F_t^{(1)}(F_0)\right)$, where
$F_t^{(1)}(F_0)$ is the solution $F_0^{(1)}$ of \eqref{eq:sabr2a},
with initial condition $F_0^{(1)}\!\!=\!\!F_0\,$,
and $\,F_t^{(2)}(F_t^{(1)}(F_0))\,$ is the solution $F_t^{(2)}$ of Equation~\eqref{eq:sabr2b},
with initial condition $F_0^{(2)}\!=\!F_t^{(1)}(F_0)$.

According to Section~\ref{sec:diff}, Equation~\eqref{eq:sabr2b} which obeys the
SABR dynamics with zero correlation, is essentially a CEV process whose solution is
\beqn
F_t^{(2)}|\sigma_t,I_0^t\sim \mathcal{CEV}_\beta\left(F_0^{(2)},(1-\rho^2)\sigma_0^2 tI_0^t\,\right)\,.
\eeqn 

This leaves us with Equation~\eqref{eq:sabr2a}, which has no known analytic solution.
Choi {em et al}\cite{CHK} suggested a frozen coefficient approximation in which
they approximated $[F_t^{(1)}]^{(1-\beta)}$ by $F_0^{(1-\beta)}$,
accordingly,

\beqn\label{eq:approx}
\begin{aligned}
\frac{dF_t^{(1)}}{F_t^{(1)}}=&\,\frac{\rho\sigma_t}{[F_t^{(1)}]^{(1-\beta)}}\,dZ_t\\[5pt]
\approx&\,\frac{\rho\sigma_t}{F_0^{(1-\beta)}}\,dZ_t\,.
\end{aligned}
\eeqn

The solution $\tilde{F}_t^{(1)}$ of Equation~\eqref{eq:approx} is

\beqn
\tilde{F}_t^{(1)}(F_0)=F_0\,\exp\left(\frac{\rho}{\nu F_0^{(1-\beta)}}(\sigma_t - \sigma_0) - 
\frac{\rho^2\sigma_0^2 t I_0^t}{2F_0^{2(1-\beta)}}\right)\,.
\eeqn

Finally, conditional on $\sigma_t$ and $I_0^t$, the CEV approximation of the general
SABR process results in a terminal forward price $F_t$ according to

\beqn
F_t|\,\sigma_t,I_0^t\sim\,\mathcal{CEV}\left(\tilde{F}_t^{(1)},\,(1-\rho^2)\sigma_0^2\, t\, I_0^t\right)
\eeqn
As is the case with almost all approximation approaches of this nature,
the accuracy of the final result increases in inverse proportion to
the time step $t$.

\section{Implied Volatility}
Given an option pricing model, the {\em implied
volatility} is a value, reflective of the variance of the stochastic process,
that, when plugged into the model, gives the observed option price.
Derivatives traders prefer implied volatility as a price quotation mechanism.
This is due to the fact that it can provide stable option quotation that
does not need to change with every  movement of the underlying.
Moreover, implied volatility reveals the nature of option prices as
a function of option strike---generally referred to as the volatility smile.

Let $B(\sigma;u,k,\tau)$ be the price of an option of strike $k$, on an underlying
of forward price $u$ and price volatility $\sigma$,
and which expires at time $\tau$.
Given a specific tuple $(v,u_0,k,\tau)$,
where $v$ the price of the option at time $t=0$,
the implied volatility of the option is
\begin{subequations}
\beqn
\sigma^{imp} = V(v;u_0,k,\tau)\,,
\eeqn
that satisfies
\beqn
B(V(v;u_0,k,\tau);u,k,\tau)=v\,,
\eeqn
\end{subequations}
\ie, $V(.)=B(.)^{-1}$.


When the dynamics of the underlying are believed to be governed by a
geometric Brownian motion, the pricing model is log-normal and the
implied volatility is often referred to as the Black
or Black-Scholes volatility, $\sigma^{BS}$.
For instance, in the case of a call option,
$\sigma^{BS}$ can be obtained from 

\beqn
%\begin{aligned}
B(\sigma;u_0,k,\tau)=u_0\,N(d_1)\,-\,k\,N(d2)\quad\quad\text{with}\quad\quad
d_{1,2}=\frac{\log(u_0/k)}{\sigma^{BS}\sqrt{\tau}}
\pm \frac{\sigma^{BS}\sqrt{\tau}}{2}\,,
%\end{aligned}
\eeqn
where $N(.)$ is the CDF of the standard Gaussian distribution.

Similarly, when the dynamics are Gaussian in nature, the 
pricing model is Bachelier and the implied volatility is known as the
normal volatility, $\sigma^{N}$, which
in the case of a call option should satisfy

\beqn
B(\sigma;u_0,k,\tau)=(u_0-k)\,N(d_N)\,+\,\sigma^{N}\sqrt{\tau}\,n(d_N)\quad\quad
\text{with}\quad\quad d_n=\frac{u_0\,-\,k}{\sigma^N\sqrt{\tau}}\,,
\eeqn
where $n(.)$ is the PDF of the standard Gaussian distribution.

More relevant to the SABR model is the notion of the CEV implied volatility,
$\sigma^{CEV}$,
which is the volatility associated with the constant elasticity of variance
model.
In the case of a call option, $\sigma^{CEV}$ obeys (see \cite{CW})

\beqn
\begin{aligned}
B(\sigma;u_0,k,\tau)=&\,\bar{F}_{\chi^2}\left(\frac{k^{2\beta_*}}{[\sigma^{CEV}\beta_*]^2\,\tau};
2+\frac{1}{\beta_*},\frac{1}{[\sigma^{CEV}\beta_*]^2\,\tau}\right)\,-\\[5pt]
&k\,
F_{\chi^2}\left(\frac{1}{[\sigma^{CEV}\beta_*]^2\,\tau};
\frac{1}{\beta_*},\frac{k^{2\beta_*}}{[\sigma^{CEV}\beta_*]^2\,\tau}\right)\,,
\end{aligned}
\eeqn
where $\beta_*\!=\!1\!-\!\beta$, $\,\bar{F}_{\chi^2}(.)\!=\!1-\!F_{\chi^2}(.)$ and
$F_{\chi^2}(x;r;x_0)$ is CDF of the non-central chi-squared distribution 
of $r$ degrees of freedom and $x_0$ non-centrality parameter.

It goes without saying that any implied volatility can be converted to
any other implied volatility.
This is done by plugging the given implied volatility into the
corresponding option pricing model.
Once the option price is obtained, it can be used to obtain the
desired implied volatility by, usually numerically, solving the
relevant inverse pricing model $B^{-1}(.)$.

\section{Asymptotic Representations}
Since deriving closed form analytic formulas for derivatives pricing is,
most of the time, impossible or impractical,
a significant amount of research effort has been directed towards
finding effective semi-analytic approximations.
These are usually formulas for expressing implied volatility in terms
of model parameters.
In this section, we focus on semi-analytic representations of the SABR model.
All formulas presented in this section will use the following
normalized quantities,

\beqn
\begin{aligned}
    f_t =&\frac{F_t}{F_0}\quad(f_0=1)&&k=\frac{K}{F_0}\\
    \hat{\sigma}_t=&\frac{\sigma_t}{\sigma_0}\quad(\hat{\sigma}_0=1)&&
    \alpha = \frac{\sigma_0}{F_0^{\beta_*}}\,,
\end{aligned}
\eeqn
where $K$ is the option strike price.
We will denote option expiry by $T$.

\subsection{Hagan {\em et al}'s (HKLW) Formula}
The Hagan approximation (2002)\cite{HKLW} is one of the most widely used
semi-analytic approximations of the SABR model.
It is very effective near the money and for small times to maturity.
Hagan {\em et all} presented $O(T)$ approximations for both the normal and the
Black-Scholes volatilities.
Their normal approximation is more accurate.
In a later paper, Hagan {\em el al} (2014)\cite{HKLW1},
derived a modified more accurate version of their earlier approximation.
In the modified Hagan model, the normal implied volatility of the SABR model
is given by the following chain of definitions:

\beqn\label{eq:hklw}
\begin{aligned}
\sigma^N=&\alpha\,\frac{q_N}{q}\,H(z)\,(1\,+\,h_N\,T),\quad\quad
z=\frac{\nu}{\alpha}\,q\\
q=&\,\int_1^kk^{-\beta}\,dk=
\begin{cases}
    \frac{k^{\beta_*}-1}{\beta_*}\quad\text{if}\quad 0\le\beta<1\\
    \log k\quad\text{if}\quad \beta=1
\end{cases}\\
q_N=&k-1\quad\quad(z_N=\frac{\nu}{\alpha}\,q_N)\\
H(z)=&\,\frac{z}{x(z)},\quad x(z)=\,\log\left(\frac{V(z)+z+\rho}{1+\rho}\right),
\quad V(x)=\sqrt{1+2\rho z+z^2}\\[5pt]
h_N=&\,\log\left(k^{\beta/2}\frac{q}{q_N}\right)\frac{\alpha^2}{q^2}\,+\,
\frac{\rho}{4}\frac{k^\beta-1}{k-1}\alpha \nu \,+\,
\frac{2-3\rho^2}{24}\nu^2\,.
\end{aligned}
\eeqn

There are three special cases concerning Equation~\eqref{eq:hklw}.
First, when $z\!=\!0\,$, $H(z)\,$ converges to $1$.
Secondly, when $\rho\!=\!0\,$, $H(z)\,$ simplifies to $H(z)\!=\!z/arcsinh(z)$.
Finally, the third case is the case of $\rho\!=\!-1\,$.
In such a case,
the value of $H(z)$ becomes
\beqn
H(z)=
\begin{cases}
    0\quad\quad\text{if}\quad z=0\\
    \infty\quad\quad\text{otherwise}\,.
\end{cases}
\eeqn

\subsection{Choi and Wu's CEV Formula}
In 2021 Choi and Wu\cite{CW} presented a powerful formula for the
approximation of $\sigma^{CEV}$, based on
Paulot's (2015) work\cite{P}.
This formula is better at handling out of the money options,
large initial volatility and large vol-of-vol.

\section{Option Portfolio}
In the preceding sections we developed a framework for simulating
the correlated movement of forward prices and volatilities at
future time horizons.
Such framework is useful for pricing options written at time $t_0$ 
with expiry that
matches any of the simulation horizon.
This assumes that the starting values of $\sigma, \nu, \rho, \beta\,$
are known at time $t_0$ and either, remain constant, or are
given as known function in time.
Moreover, options written on different tenors of the same
interest rate curve have to reflect the stochastic dependence
of the forward prices of the various tenors.
Otherwise, simulation results may be totally unrealistic or even non-arbitrage-free.
In this section we will address these two issues.

\subsection{Forward Swap Volatility}
Let us start by assuming the following risk neutral dynamics for the instantaneous forward rate $f(t,T)$,
which is the rate at time $T$ as observed at time $t$,

\beqn\label{eq:forwards}
\begin{aligned}
df(t,T)=&\,\sum_{i=1}^n\,\sigma_i(t,T)\,\bar{\sigma}_i(t,T)\,dt\,+\,\sum_{i=1}^n\,\sigma_i(t,T)\,dW_i(t)\\
\bar{\sigma}_i(t,T)=&\,\int_t^T\,\sigma_i(t,u)\,du\,,
\end{aligned}
\eeqn
where $dW_i(t)$ are $n$ standard uncorrelated Wiener increments.

Employing Ito's lemma and the arguments of HJM (see~\cite{HJM}), we get
\beqn
\frac{dP(t,T)}{P(t,T)}=\,r(t)\,dt \,-
\,\sum_{i=1}^n\,\bar{\sigma}_i(t,T)\,dW_i(t)\,,
\eeqn
where $r(t)=f(t,t)$.

The interest rate forward swap rate $S(t_x,t_s,t_e)$ is the coupon rate of the
fixed leg of the swap, which is an $\mathcal{F}_t$-martingale under the annuity measure.
The parameters $t_s\,$, $t_e\,$ are the start and end dates of the swap, respectively,
and $t_x$ is the rate set date.
Let the swap-corresponding annuity be defined as
\beqn
A(t)=\sum_{i=0}^n\,\pi_i\,P(t,T_i)\qquad t\le t_s\,,
\eeqn
where $n$ is the number of payments, $T_i$ is the $i$-th payment date,
$T_0=t_s$, $T_n=t_e$ and
$\pi_i$ is a calendar factor.
The value of $S(t_x,t_s,t_e)$, assuming simple calendar conventions, can be expressed as

\beqn
S(t_x,t_s,t_e)=\frac{1-P(t_x,t_e)}{A(t_x)}\,.
\eeqn

Let us write the risk neutral dynamics of $S(t_x,t_s,t_e)$ as
\beqn
dS(t,t_s,t_e)=\mu(t)\,dt\,+\,\sigma_{s,e,i}(t)\,dW_i(t)\,.
\eeqn

Employing Ito's lemma, we have, for $t\le t_s$
\beqn
\begin{aligned}
\sigma_{s,e,i}(t)=&\,\frac{1-P(t,t_e)}{A^2}\sum_{j=0}^n\pi_i P(t,T_j)\,\bar{\sigma}_i(t,T_j)\,+\,\frac{P(t,t_e)}{A(t)}\,\bar{\sigma}_n(t,t_e)\\
\sigma_{s,e,i}(t)=&\,\frac{1}{A(t)}\left(\frac{1-P(t,t_e)}{A(t)}\sum_{j=0}^n\pi_i P(t,T_j)\,\bar{\sigma}_i(t,T_j)\,+\,P(t,t_e)\,\bar{\sigma}_n(t,t_e)\right)\\
\sigma_{s,e,i}(t)=&\,\frac{1}{A(t)}\left(S(t,t_s,t_e)\,\sum_{j=0}^n\pi_i P(t,T_j)\,\bar{\sigma}_i(t,T_j)\,+\,P(t,t_e)\,\bar{\sigma}_n(t,t_e)\right)\,.
\end{aligned}
\eeqn

Substituting from Equation~\eqref{eq:forwards}, we get

\beqn\label{eq:fsvol}
\sigma_{s,e,i}(t)=\frac{1}{A(t)}\left(S(t,t_s,t_e)\,\sum_{j=0}^n\pi_i\, P(t,T_j)\int_t^{T_j}\sigma_i(t,u)\,du\,+\,P(t,t_e)\int_t^{t_e}\sigma_i(t,u)\,du\right)\,.
\eeqn

Equation~\eqref{eq:fsvol} represents the forward swap volatility in terms of the terms
structure volatilities of the instantaneous forward rates.
Using this equation one can construct a set of arbitrage-free swap volatilities,
given a set of instant forward volatilities.
Similarly, given a set of {\em consistent} swap volatilities, one should be able to deduce,
at least numerically, the
instantaneous forward volatilities.
The consistency condition is needed due to the fact that the nonlinear function of Equation~\eqref{eq:fsvol}
is necessarily invertible.

\subsection{Historical or Implied Volatility}
Since true volatilities are not market observable,
a proxy is always needed for their estimation.
Practitioners either, resort to the liquid options market and use the implied volatility as a proxy,
or resort to history and use realized volatility as a proxy.
We propose the use of the options market implied volatilities to estimate the current
instantaneous forward volatilities.
For risk management, a stochastic volatility model is usually needed.
The model can be in the form of a stochastic differential equation, or in the 
form of a time series GARCH style model (see Tawfik (2025)\cite{T} for an example).
The volatilities estimated from the options market can represent the starting
values of the stochastic model.
The parameters of the model can then be estimated from the historical data.
Even then, one has a choice of either, resorting to the history of swaption implied volatilities,
or the history of realized swap volatilities.

\subsection{Instantaneous Forward Volatility Model}


\subsection{Options on Different Tenors of Same Curve}
Under the SABR model, the stochastic dependence across options written on different tenors of the same curve
is made of two components.
First is the correlation structure among the forwards themselves.
Second is the correlation structure among the SABR volatility.

If $F_t^i$ and $\sigma_t^i$ are the price and the SABR volatility
of the forward contract associated with the $i$'s tenor at time t,
then the following represents the associated Wiener increments

\begin{subequations}
    \begin{align}
        dW_t^i=&\,\frac{dF_t^i}{\sigma_t^i\,(F_t^i)^\beta}\\
        dZ_t^i=&\,\frac{dZ_t^i}{\nu^i\,\sigma_t^i}\\
        [dW_t^i,\,dZ_t^i]=&\,\rho^i\,dt\,,
    \end{align}
\end{subequations}
where $\beta^i$ and $\nu^i$ are the tenor specific SABR parameters.
A time series for $dW_t^i$ and $dZ_t^i$ for each of the $N_T$ tenors can be constructed
based on the forward prices and the corresponding shortest expiry option prices.
Given a forward price, SABR model parameters are obtained by fitting the model
to the prices of options of same expiry and various strikes.
Having the individual tenor-specific time series, of size $N_t$ each,
we then construct the two associated correlation matrices of the space of all tenors:
the forward price correlation matrix $C_f$ and the volatility correlation matrix $C_\sigma$.
The correlation matrices $C^{N_t\times N_T}$ can then be diagonalized, using the singular value decomposition (SVD),
into the form $U^{N_t\times N_t}\,\Sigma^{N_t\times N_T}\,V^{N_T\times N_T}$, where $\Sigma$ is a diagonal
matrix whose elements are in a
non-decreasing order.
It is typical for the correlation structure of the whole space to be explained
be no more than the four largest values of $\Sigma$.

To include the stochastic dependence among options written on the same
curve we modify Equations~\eqref{eq:sabr} to read

\begin{subequations}\label{eq:sabr_corr}
    \begin{align}
        dF_t^i=\,&(F_{t}^j)^{\beta} \sigma_t dW_t^j\quad&\beta\in[0,1],\,F_0^i>0\label{eq:sabr_corr_a}\\
        d\sigma_t^i =\,&\nu \sigma_t^i dZ_t&\nu\ge 0,\,\sigma_0^i>0\label{eq:sabr_corr_b}\\
        [dZ_t^i,dW_t^i] =\,& \rho^i dt&\rho \in[-1,1]\\
        dZ_t^i =\,&(1-\sum_{j=1}^{N_\sigma}\gamma_j^2)^\half d\tilde{Z}_t^i + \sum_{j=1}^{N_\sigma}\gamma_{ij} d\tilde{Z}_t^j\\
        dW_t^i =\,&(1-\sum_{j=1}^{N_f}\lambda_{ij}^2)^\half d\tilde{W}_t^i + \sum_{j=1}^{N_f}\lambda_{ij} d\tilde{W}_t^j
    \end{align}
\end{subequations}
where $N_f$ and $N_\sigma$ are the number of the largest eigenvalues $\lambda_{ij}$
sufficient to explain the total variances of the set of forwards and their option volatilities,
respectively.


\subsection{Options on Same Underlying}
When given a portfolio of options written on the same underlying,
but with different strike prices and expiry dates,
the time dependency of the SABR parameters is implicitly hidden
in the prices of these options.
In this section, we will develop an approach for extracting the
implied risk-neutral time dependency information from an option
portfolio.

Let $\Pi$ be a portfolio of a collection of at-the-money options.
Each option $C_i\in\Pi$ is associated with a tuple of
SABR parameters $(\sigma,\nu,\rho,\beta,\tau)_i$.
The SABR parameters of $C_i$ are identified by examining
all available option prices of same expiry on the same underlying.
For SABR parameter identification techniques, see 
 


\section{Implied Volatility Surface}
Let $\mathbb{D}$ be defined as
\beqn
\mathbb{D}=(0,\infty)\times[0,1]\times[-1,1]\times[0,\infty)\quad\quad
x\in\mathbb{D}=(\sigma,\beta,\rho,\nu)\,,
\eeqn
where $\sigma,\beta, \rho, \nu$ are the SABR parameters
associated with an option on a forward contract
of tenor $t$ and which expires at time $\tau$.
Also, let the functions $\hat{S}, \hat{T}$ be defined according to
\beqn
\begin{aligned}
    \hat{S}:&\;I_1\times I_2\longrightarrow \,\mathbb{D}\\
    \hat{S}:&\;(i,j)\mapsto (\sigma_{ij},\beta_{i,j},\rho_{i,j},\nu_{i,j})\\
    \hat{T}:&\;I_1\times I_2\longrightarrow \;\mathbb{R}^+\!\times\mathbb{R}^+\\
    \hat{T}:&\;(i,j)\mapsto (t_i,\tau_j)\,
\end{aligned}
\eeqn
where, $I_1$ and $I_2$ are two finite ordered index sets.
Now, we define the implied volatility surface as the function
$S_h(t,\tau;\hat{T}^{-1},\hat{S}):\mathbb{R}^+\!\times\!\mathbb{R}^+\rightarrow \mathbb{D}\;$
for all $(t,\tau)\in \mathbb{R}^+\!\times\mathbb{R}^+$,
where $h$ is the valuation date, which can be in the future.
In other words, the volatility surface is a function that maps
any tenor-expiry combination into a set of SABR parameters by interpolating
the set of parameters given for a select finite set of options.
Typically, the set of given parameters forms a grid pattern, \ie
$\hat{S}$ is surjective.

In order for the volatility surface, $S_h$, to be useful in proving a price
for an option, it has to be accompanied by a price $F_h(t)$ for
the underlying forward contract of tenor $t$.
In case of the SABR model, $F_h$ and $S_h$ are simulated together due to
their coupled processes.


\section{Control Variable}
It is not uncommon in a risk management scenario for the set $I_1\times I_2$
to have a size in the hundreds.
Estimating portfolio risk requires the simulation of the volatility surface over
hundreds of forward horizons.
It is not hard to see how the computational complexity grows very rapidly.
As a result, each simulation of the volatility surface has to be
limited to a number of paths that is far less than that which
needed for accurate simulation of an option price.
To alleviate that deficiency, we resort to a technique know as the
{\em control variable} technique as a means of variance reduction and,
consequently, accuracy improvement.

Our control variable approach uses the prices of four control options
struck at different strikes, for every node on the implied volatility
grid and for every simulation horizon, to adjust the distribution
of the forward price $F_h$ of the simulated sample paths.
The philosophy behind this stems from the fact that the
true forward price distribution can be ascertained from 
the option prices of options struck at different strikes.
Given a forward price $F_0(t)$ and an implied volatility
set $S_0(t,\tau)$, one can accurately price options of various strikes
that share the tenor $t$ and the expiry $h$.
This is usually done using a semi-analytic technique, such as those
discussed in Section~\ref{sec:approx},
additional Monte Carlo simulation or a partial differential equation (PDE)
solver such as the finite difference method (FDM).
Once the option prices are obtained, we compare their values
with those predictable by our simulation.
The discrepancy, is then, corrected by proportionately
adjusting the values of $F_h$
of the simulation paths.

\section{The Algorithm}

\section{Negative Rates}

\section{Results}

\section{PCA Complexity Reduction}

\section{Bayesian Volatility}
Assuming a standard filtered probability space $(\Omega, \mathcal{F}, \mathcal{F}_t, P)$,
let the processes $F_t(F_0,\sigma_t,I_t),\,\sigma_t(\sigma_0)\,\text{and}\,I_t(\sigma_t)$
be $\mathcal{F}_t$-measurable according to the SABR dynamics, stated in the preceding sections.
In this section, we are concerned with finding the best estimate for
$\sigma_t$, where $F_t$ has been generated from an exogenous process.
Using Bayesian statistics, we can write

\beqn
P(\sigma_t=s|F_t=f)=\frac{P(F_t=x|I_t=a,\sigma_t=s)P(I_t=a|\sigma_t=s)P(\sigma_t=s)}{P(F_t=f)}\,.
\eeqn
Or more precisely,

\begin{subequations}
    \begin{align}
        P&(\sigma_t=s|F_t=f)=\,\frac{P(F_t=f|I_t\in da,\sigma_t\in ds)P(I_t=a|\sigma_t\in ds)P(\sigma_t=s)}{P(F_t=f)}\\
        P&(F_t=f)=\int_0^\infty \int_0^\infty P(F_t=f|I_t\in da,\sigma_t\in ds)\,P(I_t=a|\sigma_t\in ds)\,P(\sigma_t=s)\, da\,ds\,.
    \end{align}
\end{subequations}
Since we have already shown, in the early part of this paper, how to compute the distributions
$P(F_t=f|I_t\in da,\sigma_t\in ds),\,P(I_t=a|\sigma_t\in ds)\,\text{and}\,P(\sigma_t=s)$,
we can now, at least numerically, compute the distribution $P(\sigma_t=s|F_t=f)$ and
its CDF $Q_\sigma(s|F_t=f)$ defined as,

\beqn
\begin{aligned}
    Q_\sigma(s|F_t=f)=&\,P(\sigma_t\le s|F_t=f)\\
        =&\,\int_0^s P(\sigma_t=s|F_t=f)\,ds\,.
\end{aligned}
\eeqn
To sample the distribution $Q_\sigma(.)$, we need to compute $Q_\sigma^{-1}(U)$,
where $U\sim Uniform(0,1)$ and $Q_\sigma^{-1}$ is the inverse of $Q_\sigma(.)$




\bigskip
\begin{thebibliography}{99}
    \bibitem{AW}
        Abate, Joseph and Whitt, Ward, 2006, "A Unified Framework for Numerically Inverting Laplace Transforms".
        doi:10.1287/ijoc.1050.0137
    \bibitem{CHK}
        Choi, Jaehyuk; Hu, Lilian and Kwok, Yue Kuen, 2025, "Efficient and Accurate Simulation of the Stochastic-alpha-beta-rho Model,"
        arXiv:2408.01898v2
    \bibitem{CKTW}
        Choi, Jaehyuk; Kwak, Minsuk; Tee, Chyng Wen; and Wang, Yumeng, 2022, "A Black-Scholes user's guide to the
        Bachelier model," {\em Journal of Futures Markets} 42(5):959-980. https://ink.library.smu.edu.sg/lkcsb\_research/6869
    \bibitem{CSC}
        Cai, Ning; Song, Yingda and Chen, Nan, 2017, "Exact Simulation of the SABR Model,"
        {\em Operations Research} 65(4):931-951. doi:10.1287/opre.2017.1617
    \bibitem{CW}
        Choi, Jaehyuk and Wu, Lixin, 2020, "The Equivalent Constant-elasticity-of-variance (CEV) Volatility
        of the Stochastic-Alpha-Beta-Rho (SABR) Model,"
        arXiv:1911.13123v3
    \bibitem{HKLW}
        Hagan, Patrick; Kumar, Deep; Lesnieweski, Andrew S. and Woodward, Diana, 2002,
        "Managing Smile Risk," {\em Wilmott} 1(2002):84-108.
    \bibitem{HKLW1}
        Hagan, Patrick; Kumar, Deep; Lesnieweski, Andrew S. and Woodward, Diana, 2014,
        "Arbitrage-Free SABR," {\em Wilmott} (2014):60-75. doi:10.1002/wilm.10290
    \bibitem{HJM}
        Heath, D., Jarrow, R. and Morton, A. (1992), "Bond Pricing and the Term Structure of Interest Rates:
        A New Methodology for Contingent Claims Valuation," {\em Econometrica} 60(1):77-105.
        doi:10.2307/2951677
    \bibitem{I}
        Islah, Othmane, 2009, “Solving SABR in Exact Form and Unifying It with LIBOR Market Model,”
        SSRN: http://papers.ssrn.com/sol3/papers.cfm?astract-id=1489428
    \bibitem{K}
        Kang, C., 2014, "Simulation of the Shifted Poisson Distribution with an Application
        to the CEV Model," {Management Science and Financial Engineering}, 20:27-32.
        doi:10.7737/MSFE.2014.20.1.027
    \bibitem{L}
        Cheng, Luo, 2016, "On the Calibration of the SABR Model and its
        Extensions," Thesis submitted as part of the requirements for the award
        of the MSc in Mathematics and Finance, Imperial College London. 
    \bibitem{MG}
        Makarov, R.N. and Glew, D., 2010, "Exact Simulation of Bessel Diffusions,"
        {\em Monte Carlo Methods and Applications} 16:312-347.
        doi:10.1515/mcma.2010.010. arXiv:0910.4177
    \bibitem{MY}
        Matsumoto, Hiroyuki and Yor, Marc, 2005, "Exponential functionals of Brownian motion, I: Probability
        laws at fixed time," {\em Probability Surveys} 2:312-347.
        doi:10.1214/154957805100000159 
    \bibitem{P}
        Paulot, Louis, 2015, "Asymptotic Implied Volatility at the Second Order with
        Application to the SABR Model," In: Friz, P., Gatheral, J., Gulisashvili,
        A., Jacquier, A., Teichmann, J. (eds) {\em
        Large Deviations and Asymptotic Methods in Finance}.
        Springer Proceedings in Mathematics \& Statistics, (110)
        Springer, Cham. https://doi.org/10.1007/978-3-319-11605-1\_2
        arXiv:0906.0658v3
    \bibitem{SGO}
        van der Stoep, Anthonie and Grzelak, Lech Aleksander and
        Oosterlee, Cornelis W., 2015, "The Time-Dependent FX-SABR Model:
        Efficient Calibration Based on Effective Parameters,"
        {\em International Journal of Theoretical and Applied Finance}, 18(6),
        SSRN: https://ssrn.com/abstract=2503891 
        doi:10.2139/ssrn.2503891
    \bibitem{T}
         Tawfik, Maged, 2025, "Monte Carlo Simulation of SOFR Yield Curve," {\em Working Paper}.
\end{thebibliography}
\end{document}
